
\documentclass{article}
\usepackage{graphicx} % Required for inserting images

\usepackage{amsmath}
\usepackage{amsfonts}
\usepackage{tikz-cd}

\title{test}
\date{May 2025}

\begin{document}

\vspace{0.5cm}

\textbf{Exercise 7.1}

\vspace{0.5cm}

First we need a lemma, if we have two normal matrices $A,B$ where $A^*B =BA^*$ and $AB = BA$ then they can be simultaneously diagonalized. All we need is that if $V_\lambda$ is an eigenspace for $A$ that $B(V_\lambda) \subseteq V_\lambda$ and $B^*(V_\lambda) \subseteq V_\lambda$ since then $B$ is normal when its range and domain is restriced to $V_\lambda$ and it can be diagonalized on $V_\lambda$ so there is a basis of simultaneous eigenvectors for $A$ and $B$. 

\vspace{0.3cm}

Now to the solution of the exercise. If $A \in\mathcal A$ in our abelian subgroup of $U(n)$ then the $V^{(A)}_\lambda$ decompose the entire $\mathbb C^n$ and if this decomposition doesnt diagonalize all matrices in $\mathcal A$ then pick one $A_1$ that isnt diagonal on these subspaces and further decompose each $V^{(A)}_\lambda$ into eigenspaces for $A_1$ (can be done by our lemma) and if this decomposition of $\mathbb C^n$ still doesnt diagonalize all matrices in $\mathcal A$ then continue this process of further decomposing, this process has to eventually stabilize after a maximum of $n$ steps since at that point either all matrices are diagonal on each subspace OR each subspace has dimension $1$ (since each step of this process strictly reduces the dimension of at least 1 subspace in our decomposition which means after n steps they are all dimension 1) and of course that means each matrix is diagonalized so we are done.

\vspace{0.5cm}

\textbf{Exercise 7.7}

\vspace{0.5cm}

If $e_i$ is an orthonormal basis for $V_\pi$ then $\alpha_i := (-,e_i)$ are an orthonormal basis for $V_{\pi^*}$ then

\begin{align*}
   \chi_{\pi^*}(x) &= \sum\limits_i (\pi^*(x)\alpha_i, \alpha_i) \\
   & = \sum\limits_i (v_{\alpha_i}, v_{\pi^*(x)\alpha_i}) \\
   & = \sum\limits_i (e_i, \pi(x)e_i) \\
   & = \overline{\sum\limits_i (\pi(x)e_i, e_i)} = \overline{\chi_\pi(x)}
\end{align*}

\vspace{0.5cm}

\textbf{Exercise 7.8} (incomplete)

\vspace{0.5cm}

$(a) \to (b)$

\vspace{0.3cm}

\noindent Compose the linear isomorphism of representations $V_\pi \rightarrow V_{\pi^*}$ with the canonical antilinear map (described above theorem 7.2.1) $V_{\pi^*} \rightarrow V_\pi$ and you have an antilinear bijective map $V_{\pi} \rightarrow V_{\pi}$.

\vspace{0.3cm}

$(b) \to (c)$

\vspace{0.3cm}

\noindent We can also assume that $(x,y) = (Cy,Cx)$ since the $C$ that was constructed before satisfies this. 

\vspace{0.5cm}

\textbf{Exercise 7.10}

\vspace{0.5cm}

Let $M_\pi$ be the subspace of $L^2(K)$ spanned by the matrix coefficients of the finite dimensional irreducible representation $\pi$. Let $f \in L^2(K)$ be conjugation invariant, then the projection of $f$ to the subspace $M_\pi$ is conjugation invariant aswell since $M_\pi$ is invariant under $L_k$ and $R_k$ as noted in Theorem 7.2.1. So it suffices to prove that if $f \in M_\pi$ and $(k.f)(x) = f(x)$ that $f(x) = c \text{ tr}(\pi(x))$ for some $c \in \mathbb C$.

Fix some orthonormal basis $e_i \in V_\pi$, then w.r.t this basis $f(x) = \sum\limits_{i,j}\lambda_{ij}\pi_{ij}(x)$ for some $\lambda_{ij}$ by definition of $M_\pi$. Now lets investigate each $\lambda_{ij}$, since $\pi_{ij}$ are orthogonal in $L^2(K)$ we can pull out this $\lambda_{ij}$ like so:

\begin{align*}
    \lambda_{ij} =\frac{1}{\dim(V_\pi)}  (f,\pi_{ij}) & =\frac{1}{\dim(V_\pi)} \int f(x) \overline{(\pi(x) e_i, e_j)} \ dx \\
    & =\frac{1}{\dim(V_\pi)} \overline{\int (\overline{f(x)}\pi(x) e_i, e_j)} \ dx \\
    & = \frac{1}{\dim(V_\pi)} \overline{(\pi(\bar f)e_i, e_j)} 
\end{align*}

Now since $\pi(L_kg) = \pi(k)\pi(g)$ and $\pi(R_kg) = \pi(g)\pi(k^{-1})$ for $g \in L^1$ then the condition $k.g = g$ implies $\pi(k)\pi(g) \pi(k^{-1}) = \pi(g)$ i.e that $\pi(g)$ is an intertwining operator and since $\pi$ is irreducible that $\pi(g) = \lambda  \text{Id}$. Clearly if $k.f = f$ then $k.\bar f = \bar f$. Thus $\pi(\bar f) = \lambda \text{Id}$ and we obtain:

\begin{align*}
     \lambda_{ij} = \frac{1}{\dim(V_\pi)} \overline{(\pi(\bar f)e_i, e_j)} = \frac{\bar \lambda}{\dim(V_\pi)} &= c \ \text{ (if } i = j) \\
     & = 0 \ \text{ (if } i \neq j)
\end{align*}

and thus $f(x) = c \sum\limits_{i} \pi_{ii}(x) = c \text{ tr}(\pi(x))$.

\vspace{0.5cm}

\textbf{Exercise 7.12}

\vspace{0.5cm}

If $\forall k\exists l: lkl^{-1} = k^{-1}$ then 

\begin{align*}
    \overline{\chi_\pi(k)} & =\overline{\text{tr}(\pi(k))}  \\
    & = \text{tr}(\pi(k)^*) \\
    & = \text{tr}(\pi(k^{-1}))\\
    & = \text{tr}(\pi(l)\pi(k)\pi(l)^{-1})) \\
    & = \text{tr}(\pi(k)) = \chi_\pi(k)
\end{align*}

and $\chi_\pi$ is real. 

Now the other direction, we will assume for a contradiction that 
$$\exists k \forall l : lkl^{-1} \neq k^{-1} \ \text{ and all } \ \chi_\pi \  \text{ are real}$$ 
Now adopting the same notation as in exercise 7.10 and assuming the result of that exercise, if $g \in L^2(K)$ and $x.g = g$ for all $x \in K$ then we can according to exercise 7.10 write 
$$g(x) = \sum\limits_\pi c_\pi \chi_\pi(x)$$      
where the sum converges in $L^2$, and since $f(x) \mapsto f(x^{-1})$ mapping a function $f(x)$ to $f(x^{-1})$ is an isometry in $L^2(K)$, which means its a bounded operator and commutes with infinite sums, then 
$$g(x^{-1}) = \sum\limits_\pi c_\pi \chi_\pi(x^{-1}) = \sum\limits_\pi c_\pi \overline{\chi_\pi(x)} = \sum\limits_\pi c_\pi \chi_\pi(x) = g(x)$$
so that necessarily $g(x) = g(x^{-1})$ if $\forall x \in K :x.g = g$. Now let's construct a function $g$ with $x.g = g$ for all $x \in K$ but $g(x) \neq g(x^{-1})$ which will be our desired contradiction. Remember we assumed $\exists k \forall l : lkl^{-1} \neq k^{-1}$ then let 
$$E = \{ lkl^{-1}: l \in K\}$$ 
It's a compact set and 
$$k^{-1} \not \in E$$ 
so there exists some continuous $f$ with $f = 1$ on $E$ and $f=0$ on some neighborhood of $k^{-1}$ and $0 \leq f \leq 1$ everywhere. Then consider 
$$g(x) = \int f(yxy^{-1}) \ dy$$
$g(k) = 1$ since $f = 1$ on $E$ and we chose the normalization $\mu(K) = 1$, but $g(k^{-1}) < 1$ since for all $y$ in some neighborhood $U$ of $y = 1_K$ we have $f(yx^{-1}y^{-1}) = 0$ and thus

$$\int f(yx^{-1}y^{-1}) \ dy = \int_{K - U} f(yx^{-1}y^{-1}) \ dy \leq 1 \cdot (1 - \mu(U)) < 1$$

since $\mu(U) > 0$ for all nonempty open sets $U$ (corollary 1.3.6)!

\end{document}
